\documentclass[11pt,a4paper]{article}

\usepackage[T2A]{fontenc}
\usepackage[utf8]{inputenc}
\usepackage[russian]{babel}
\usepackage[a4paper,hmargin=1.5cm,vmargin=1.5cm]{geometry}
\usepackage{amsmath,amssymb}
\usepackage{array,booktabs}
\usepackage{eso-pic}

\setlength{\parindent}{0pt}
\setlength{\parskip}{4pt}
\pagestyle{empty}

% Карточка варианта: номер, функции для вронскиана, уравнение, частное решение
% Размеры карточки: сетка 2x2 на всю страницу с зазором \cardgap; рамок нет, только две линии разреза по центру
\newlength{\cardgap}\setlength{\cardgap}{0.4cm}
\newlength{\cardw}\newlength{\cardh}
\newcommand{\card}[4]{%
  {\fboxrule=0pt\fbox{\begin{minipage}[t][\cardh][t]{\cardw}
    \vspace{4pt}
    \textbf{Вариант #1}

    \smallskip
    Фамилия: \rule{5cm}{0.4pt}

    \medskip
    \textbf{1.} \ Найти вронскиан функций $#2$. Линейно зависимы ли они?

    \medskip
    \textbf{2.} \ Найти общее решение уравнения
    \[
      #3\,,
    \]
    если известно частное решение $y_1 = #4$.
  \end{minipage}}}}

\begin{document}

% Страница с карточками: поля 5 мм (меньше принтер обычно не пропечатывает)
\newgeometry{margin=0.5cm}
\AddToShipoutPictureBG*{%
  \AtPageCenter{\makebox(0,0){\rule{0.4pt}{\dimexpr\paperheight-1cm\relax}}}%
  \AtPageCenter{\makebox(0,0){\rule{\dimexpr\paperwidth-1cm\relax}{0.4pt}}}}
\setlength{\cardw}{\dimexpr(\textwidth-\cardgap)/2-2\fboxsep-2\fboxrule\relax}
\setlength{\cardh}{\dimexpr(\textheight-\cardgap)/2-2\fboxsep-2\fboxrule-2pt\relax}

% ---------------------------------------------------------------------
% Страница 1: варианты для печати и разрезания
{\setlength{\parskip}{0pt}\offinterlineskip
\noindent\card{1}{e^{2x},\ x e^{2x}}{y'' - 2xy' + (x^2 - 1)\, y = 0}{e^{x^2/2}}\hfill
\card{2}{\cos 3x,\ \sin 3x}{x y'' + 2y' + x y = 0}{\dfrac{\sin x}{x}}\par

\vspace{\cardgap}

\noindent\card{3}{x,\ x^3}{x^2 y'' - x(x + 2)\, y' + (x + 2)\, y = 0}{x}\hfill
\card{4}{e^{x},\ e^{-3x}}{x y'' - y' - 4x^3 y = 0}{e^{x^2}}\par}

\newpage
\restoregeometry
% ---------------------------------------------------------------------
% Страница 2: ответы для преподавателя
\begin{center}
{\Large\bfseries Квиз: ответы}
\end{center}

Проверяет домашнее задание к \texttt{lesson\_02.tex}: вронскиан и линейная независимость, общее решение по известному частному через формулу Лиувилля "--- Остроградского.

Во всех вариантах функции задачи 1 линейно независимы. В задаче 2 интеграл берётся без интегрирования по частям. Типы задач по вариантам:
\begin{center}
\renewcommand{\arraystretch}{1.3}
\begin{tabular}{@{}lll@{}}
\toprule
Вариант & Задача 1 & Задача 2\\
\midrule
1 & кратный корень, $W = e^{4x}$ & $p = -2x$, \ $W = e^{x^2}$ \\
2 & комплексные корни, $W = 3$ & $p = \frac{2}{x}$, \ $W = \frac{1}{x^2}$ \\
3 & степени, $W = 2x^3$ & $p = -1 - \frac{2}{x}$, \ $W = x^2 e^{x}$ \\
4 & разные корни, $W = -4e^{-2x}$ & $p = -\frac1x$, \ $W = x$ \\
\bottomrule
\end{tabular}
\end{center}

\par\medskip\textbf{Вариант 1}\par

\textbf{1.} $\displaystyle W = \begin{vmatrix} e^{2x} & x e^{2x} \\ 2e^{2x} & (1 + 2x)e^{2x} \end{vmatrix} = e^{4x}\bigl((1 + 2x) - 2x\bigr) = e^{4x} \ne 0$ \quad$\Rightarrow$\quad независимы

\textbf{2.} $p = -2x$, \ $W = e^{\int 2x\,dx} = e^{x^2}$, \ $\displaystyle \left(\frac{y_2}{y_1}\right)' = \frac{e^{x^2}}{e^{x^2}} = 1$, \ $y_2 = x e^{x^2/2}$ \quad$\Rightarrow$\quad $\displaystyle y = (C_1 + C_2 x)\, e^{x^2/2}$

\par\medskip\textbf{Вариант 2}\par

\textbf{1.} $\displaystyle W = \begin{vmatrix} \cos 3x & \sin 3x \\ -3\sin 3x & 3\cos 3x \end{vmatrix} = 3\cos^2 3x + 3\sin^2 3x = 3 \ne 0$ \quad$\Rightarrow$\quad независимы

\textbf{2.} Делим на $x$: $p = \dfrac{2}{x}$, \ $W = e^{-2\ln|x|} = \dfrac{1}{x^2}$, \ $\displaystyle \left(\frac{y_2}{y_1}\right)' = \frac{1}{x^2} \cdot \frac{x^2}{\sin^2 x} = \frac{1}{\sin^2 x}$, \ $\dfrac{y_2}{y_1} = -\operatorname{ctg} x$, \ $y_2 = -\dfrac{\cos x}{x}$

\quad$\Rightarrow$\quad $\displaystyle y = \frac{C_1 \sin x + C_2 \cos x}{x}$

\par\medskip\textbf{Вариант 3}\par

\textbf{1.} $\displaystyle W = \begin{vmatrix} x & x^3 \\ 1 & 3x^2 \end{vmatrix} = 3x^3 - x^3 = 2x^3 \not\equiv 0$, \ например, $W(1) = 2$ \quad$\Rightarrow$\quad независимы

\textbf{2.} Делим на $x^2$: $p = -\dfrac{x + 2}{x} = -1 - \dfrac{2}{x}$, \ $W = e^{x + 2\ln x} = x^2 e^{x}$, \ $\displaystyle \left(\frac{y_2}{y_1}\right)' = \frac{x^2 e^{x}}{x^2} = e^{x}$, \ $y_2 = x e^{x}$

\quad$\Rightarrow$\quad $\displaystyle y = x\,(C_1 + C_2 e^{x})$

\par\medskip\textbf{Вариант 4}\par

\textbf{1.} $\displaystyle W = \begin{vmatrix} e^{x} & e^{-3x} \\ e^{x} & -3e^{-3x} \end{vmatrix} = -3e^{-2x} - e^{-2x} = -4e^{-2x} \ne 0$ \quad$\Rightarrow$\quad независимы

\textbf{2.} Делим на $x$: $p = -\dfrac1x$, \ $W = e^{\ln x} = x$, \ $\displaystyle \left(\frac{y_2}{y_1}\right)' = \frac{x}{e^{2x^2}} = x e^{-2x^2}$, \ $\dfrac{y_2}{y_1} = -\dfrac14 e^{-2x^2}$, \ $y_2 = -\dfrac14 e^{-x^2}$

\quad$\Rightarrow$\quad $\displaystyle y = C_1 e^{x^2} + C_2 e^{-x^2}$

\end{document}
